\section{Discussion}
\label{sec:discussion}

The results support a layered interpretation of dependency structure in the Brazilian equity market, consistent with the multi-scale perspective discussed by Raddant and Di Matteo~\cite{raddant_dimatteo_2023}. At the broadest scale, the market exhibits a strong collective component. Positive average cross-correlation, periods of increased market-wide synchronization and the dominant leading eigenvalue all point to a common mode affecting a large part of the equity universe. This collective component is important because it can dominate raw correlation-based representations and obscure weaker relationships operating at more localized scales.

Beyond this market-wide component, the results reveal additional structure that cannot be reduced to simple sector classifications. Within-sector correlations are systematically stronger than between-sector correlations, while the non-leading significant eigenvectors contain recognizable but heterogeneous contrasts involving commodity-related firms, financial institutions, utilities and issuer-specific effects. The spectral results therefore suggest that intermediate-scale organization exists, but that it does not necessarily coincide with predefined economic sectors. This distinction is important when interpreting financial dependencies: statistically significant collective modes may reflect combinations of economic, institutional and issuer-specific relationships rather than isolated sector factors.

The hierarchical and network representations provide complementary evidence for this interpretation. After removal of the dominant market mode, the remaining dependency structure becomes weaker in absolute correlation terms but does not disappear. Ordered heatmaps and hierarchical clustering continue to reveal localized organization, while the lower cophenetic correlation of the group-mode representation suggests that this residual geometry is less purely tree-like. The MST and PMFG therefore add information by retaining different aspects of the dependency structure beyond a strictly hierarchical representation.

The PMFG results are particularly informative in this respect. Market-mode removal substantially reduces average edge strength, yet local clustering increases in the group-mode network. This combination indicates that weaker pairwise dependencies can still form relatively concentrated local structures once broad market synchronization is removed. The result is consistent with the role of planar filtering as a way to preserve richer local organization than the MST alone \cite{tumminello_2005_tool,tumminello_2010_correlation}. Hub reorganization further reinforces this interpretation: assets that are central in the original network are not necessarily the same assets that become central after removal of the collective component. Network centrality should therefore be understood as conditional on the dependency representation from which the graph is constructed.

The subsector-level network provides an economic summary of this residual organization. Differences between degree, weighted degree and betweenness show that connectivity, dependency intensity and bridge-like position describe distinct aspects of market structure. At the same time, the robustness results caution against interpreting these configurations as permanent characteristics. The dominant spectral component remains identifiable after collapsing multiple share classes to the issuer level, and the leading empirical eigenvalues remain distinguishable from the circular-shift surrogate benchmark. However, weaker modes are more sensitive to issuer representation, while individual network edges, hubs and communities show limited persistence across time windows. The full-sample networks should therefore be interpreted as long-run structural summaries rather than invariant market architectures.

The validation-selected Ridge comparisons show higher test QLIKE for Set B and Set C than for Set A at both horizons; date-level HAC tests favor Set A, including after Holm adjustment across the four structural contrasts. This result is specific to the selected Ridge comparisons and should not be generalized to every model or to a point-in-time deployment. HAR-RV is competitive: its paired comparison with Ridge Set A is not statistically distinguishable at the 5% level. Five-seed sensitivity shows greater QLIKE dispersion for MLP than CNN, so single-seed neural rankings are uncertain. Structural descriptors remain ex-post because they were computed from the full sample.

Overall, the study integrates spectral filtering, hierarchical organization and financial networks to characterize B3 equity dependencies, then evaluates whether adding these descriptors improves volatility forecasts. The Ridge loss comparisons indicate worse test QLIKE after adding market/RMT or network descriptors under this design; HAR and seed-sensitivity results provide relevant context. These results do not establish causal effects or real-time forecasting performance.

The findings should be interpreted in light of several limitations. First, the RMT and network descriptors used in the volatility assessment are estimated from the full synchronized sample and are therefore unavailable at historical forecast origins. The forecasting experiment should consequently be interpreted as a retrospective structural assessment rather than as a strict point-in-time test. The most important extension is to reconstruct the spectral and network descriptors recursively using rolling or expanding windows.

Second, the volatility assessment is restricted to PETR4, VALE3 and BBDC4. Although these assets provide liquid and economically distinct forecasting targets, they do not represent the full B3 cross-section. Extending the experiment to a broader set of assets would provide a stronger assessment of whether the observed relationships generalize across sectors, liquidity profiles and market conditions.

Third, long historical equity samples remain sensitive to data-quality issues associated with corporate actions, ticker transitions, share-class changes and adjustment procedures. Adjusted prices reduce these problems but do not eliminate them entirely. The issuer-level robustness analysis also shows that multiple listed share classes can materially affect some localized spectral components.

Fourth, sector and subsector classifications provide a useful economic reference but remain descriptive labels rather than complete representations of firm exposure. Companies may share dependencies through macroeconomic, liquidity or balance-sheet channels that do not align with predefined classifications. The network and eigenvector results should therefore not be interpreted as causal mappings of economic transmission mechanisms.

Finally, the 1998--2025 local extraction has not been redistributed, and its source-specific license was not verified here; the public BovDBv2 deposit ends in July 2024. Exact reproduction of the local 2025 extension requires authorized data access. The present analysis also relies on static full-sample structural descriptors and relatively simple machine-learning and neural architectures. The study does not establish causal effects of network position on future volatility, nor does it evaluate whether more sophisticated temporal or graph-based models would benefit from dynamically reconstructed network information. Future work should combine point-in-time RMT and network estimation with broader target coverage, alternative rolling-window specifications and models capable of jointly representing temporal and cross-sectional market dynamics.
